\chapter{Related Work}
\label{ch:related_work}

This chapter provides a comprehensive analysis of previous research
and state-of-the-art technologies addressing multi-language static
code analysis, structural extraction, and name resolution. The
analysis is structured to evaluate the most recent and effective
techniques, critically examining their architectural choices,
underlying mechanisms, and current limitations. Understanding these
baseline approaches is essential for positioning the framework
proposed in this thesis, highlighting how it overcomes existing
bottlenecks in building language-agnostic dependency graphs.

The chapter is divided into four main sections: approaches based on
intermediate syntactic representations, formal models for symbol and
name resolution, comprehensive multi-language static analysis
frameworks, and finally, a critical review of the direct precursors
to this work developed within our research group.

\section{Language-Agnostic Abstract Syntax Trees}

A major challenge in cross-language static analysis is the profound
heterogeneity of the Abstract Syntax Trees (ASTs) generated by
different parsers. Since standardizing the parsers themselves is
generally unfeasible, several works have attempted to unify the
resulting tree representations to allow a single backend tool to
analyze software regardless of the source language.

\subsection{Syntactic Wrapping: srcML}

\textbf{srcML} \cite{collard_srcml_2016} is an infrastructure and
XML-based format designed to represent source code by wrapping the
original raw text with syntactic XML tags.

\textit{How it works:} Instead of transforming code into an abstract
data structure that discards formatting, srcML parses the code and
produces an XML document that preserves the exact original text,
including whitespace and comments. It overlays this text with
syntactic boundaries (e.g., \texttt{<class>}, \texttt{<function>},
\texttt{<expr>}). Analysts can then use standard XML querying tools,
such as XPath, to extract facts and relationships from the codebase.

For example, given the simple C++ snippet in Listing
\ref{lst:srcml_cpp}, srcML translates it into the XML format shown in
Listing \ref{lst:srcml_xml}, preserving the original tokens while
exposing the syntax tree structure \cite{srcml_about}.

\begin{lstlisting}[language=C++,caption={Example of a simple C++ function to be analyzed by srcML.},label={lst:srcml_cpp}]
#include "rotate.h"

// rotate three values
void rotate(int& n1, int& n2, int& n3) {
  // copy original values
  int tn1 = n1, tn2 = n2, tn3 = n3;

  // move
  n1 = tn3;
  n2 = tn1;
  n3 = tn2;
}
\end{lstlisting}

\begin{lstlisting}[basicstyle=\ttfamily\footnotesize, language=XML,caption={The resulting srcML representation, keeping formatting intact but injecting syntactic tags.},label={lst:srcml_xml}]
<cpp:include>#<cpp:directive>include</cpp:directive> <cpp:file>"rotate.h"</cpp:file></cpp:include>

<comment type="line">// rotate three values</comment>
<function><type><name>void</name></type> <name>rotate</name><parameter_list>(<parameter><decl><type><name>int</name><modifier>&amp;</modifier></type> <name>n1</name></decl></parameter>, <parameter><decl><type><name>int</name><modifier>&amp;</modifier></type> <name>n2</name></decl></parameter>, <parameter><decl><type><name>int</name><modifier>&amp;</modifier></type> <name>n3</name></decl></parameter>)</parameter_list> <block>{<block_content>

  <comment type="line">// copy original values</comment>
  <decl_stmt><decl><type><name>int</name></type> <name>tn1</name> <init>= <expr><name>n1</name></expr></init></decl>, <decl><type ref="prev" /><name>tn2</name> <init>= <expr><name>n2</name></expr></init></decl>, <decl><type ref="prev" /><name>tn3</name> <init>= <expr><name>n3</name></expr></init></decl>;</decl_stmt>

  <comment type="line">// move</comment>
  <expr_stmt><expr><name>n1</name> <operator>=</operator> <name>tn3</name></expr>;</expr_stmt>
  <expr_stmt><expr><name>n2</name> <operator>=</operator> <name>tn1</name></expr>;</expr_stmt>
  <expr_stmt><expr><name>n3</name> <operator>=</operator> <name>tn2</name></expr>;</expr_stmt>
</block_content>}</block></function>
\end{lstlisting}

\textit{Weaknesses and Limitations:} While srcML is highly robust for
lightweight fact extraction and refactoring tasks, it is essentially
a syntactic markup rather than a semantic model. It does not provide
built-in name resolution or cross-file dependency linking; it simply
unifies the syntax. Furthermore, its parser ecosystem is currently
less mature than modern alternatives like Tree-sitter for supporting
the deepest semantic intricacies of all modern programming languages.

\textit{Relevance:} As the tool matures, it represents a valid future
alternative for the extraction layer. An XML-based robust parser
could potentially replace the Tree-sitter extraction phase of our
pipeline, standardizing the input format even further without
requiring the complex tree-cursor navigation and heuristic-based
identification currently employed.

\subsection{Semantic Translation: Component Identification}

Recent research by Curtis (2022) \cite{curtis_language-agnostic_2022}
and Schiewe et al.
(2022) \cite{schiewe_advancing_2022} focuses on transforming source code
into a fully language-agnostic AST, particularly in the context of
analyzing heterogeneous, cloud-native microservices. In distributed
architectures, individual microservices are often written in
different programming languages, making traditional single-language
static analysis tools ineffective for tracing system-wide
architectural dependencies (e.g., a Python service calling a Java backend).

\textit{How it works:} To overcome the language barrier, these
approaches introduce a dedicated translation pipeline. First,
language-specific parsers generate standard ASTs. Then, custom
translators map these low-level, language-specific syntax nodes into
a serialized, language-agnostic AST format (often using JSON).
Finally, generalized component parsers traverse this unified tree to
identify high-level architectural abstraction (such as REST
  endpoints, controllers, routing configurations, and data
entities) regardless of whether they were originally written in C++ or Java.

\textit{Weaknesses and Limitations:} While successful for high-level
component identification, this methodology presents two major
limitations. First, building and maintaining the translation layer is
highly expensive: every new supported language requires a custom tool
to map its AST to the universal AST. Second, translating an entire
AST into a common, standardized format is inherently a \textit{lossy}
process. To create a universal tree, one must often reduce constructs
to their lowest common denominator. Forcing completely different
semantic paradigms (e.g., Rust's trait-based composition versus
Java's object-oriented inheritance) into a generic unified node
discards subtle but crucial language-specific details. These details,
such as scoping rules and exact type bounds, are strictly required
for accurate, fine-grained name resolution.

\textit{Relevance:} Our work directly addresses these limitations by
actively avoiding the intermediate ``Universal AST'' translation step.
Instead of transforming the \textit{data} to fit a generic model, we
abstract the \textit{extraction logic}. We perform a
\textbf{heuristic-based extraction} directly on the original Concrete
Syntax Tree (CST) provided by Tree-sitter. Rather than translating the
tree, our framework considers programming languages as a combination
of independent behavioral features. By querying the native CST using
naming heuristics and checking the specific feature set of the
analyzed language (e.g., whether it supports nested modules or
package declarations), we selectively extract only the relevant
architectural elements into our Intermediate Representation
($\mathcal{D}_{IR}$).
This approach skips the heavy translation pipeline entirely and retains
the precise semantic context necessary for accurate dependency resolution,
scaling easily to new languages without writing new AST translators.

\section{Models for Name Resolution}

Once the structural components of a codebase are identified, the
subsequent challenge is resolving identifiers (e.g., variable usages,
method calls) to their concrete definitions across complex, nested scopes.

\textbf{Scope Graphs} \cite{zwaan_scope_2023}, developed by
researchers at TU Delft, provide a foundational, language-parametric
theoretical framework for name binding and resolution.

\textit{How it works:} The model abstracts away the AST entirely,
representing the program's lexical structure as a graph. Lexical
scopes are represented as nodes, and visibility or reachability rules
(such as lexical nesting or module imports) are represented as
directed edges. Resolving a name becomes a formal path-finding
problem within this graph, governed by specific regular expressions
that dictate valid resolution paths.

\textit{Weaknesses and Limitations:} The main weakness of standard
Scope Graphs is their lack of scalability. The model typically
requires a whole-program analysis, meaning the entire project must be
parsed and translated into a single, massive graph in memory before
resolution can begin. They struggle with incremental updates, making
them computationally prohibitive for very large codebases.

\textbf{Stack Graphs} \cite{creager_stack_2023} are the direct
successor to Scope Graphs, developed and heavily utilized by GitHub
to provide code navigation (e.g., "jump to definition") at an industrial scale.

\textit{How it works:} They extend Scope Graphs by introducing a
"Symbol Stack" and two new node types: \textit{push} and \textit{pop}
nodes. This crucial addition allows the graph to be built
incrementally, file by file. When a reference is queried, the
algorithm walks the graph, pushing identifiers onto the stack at
\textit{push} nodes and popping them at \textit{pop} nodes. A path is
considered valid, and the reference resolved, if the symbol stack is
completely empty when a definition node is reached.

\textit{Weaknesses and Limitations:} While theoretically elegant and
excellent for localized IDE queries, Stack Graphs possess significant
practical limitations when repurposed for total dependency
extraction. First, constructing a Stack Graph requires the creation
of highly complex and verbose grammars (using the Tree-sitter Graph
or TSG language) for every single language. Second, because
resolution in Stack Graphs relies on a pushdown automaton traversing
graph edges, finding all dependencies requires executing this
graph-walking algorithm from every single reference. In highly
interconnected codebases with deep import chains, this mechanism
suffers from a \textit{path-explosion} problem: the algorithm
evaluates thousands of combinatorial paths simultaneously to verify
stack emptiness, leading to severe memory and performance
bottlenecks. Conversely, our approach avoids this combinatorial
explosion by replacing unguided graph traversal with deterministic
algebraic lookups on a Scope Tree.

\section{Multi-Language Static Analysis Frameworks}

While the aforementioned tools focus on specific phases of analysis
(parsing or name resolution), other solutions attempt to provide
end-to-end multi-language analysis frameworks.

\textbf{Joern} \cite{joern_docs} is a prominent, open-source
platform originally developed by Yamaguchi et al. for robust,
multi-language static analysis, primarily targeting vulnerability
discovery\cite{yamaguchi_modeling_2014}.

\textit{How it works:} Joern generates a \textit{Code Property Graph
(CPG)}, a highly complex data structure that unifies the AST, the
Control Flow Graph (CFG), and the Program Dependence Graph (PDG) into
a single, queryable graph database. Analysts can interrogate this
massive graph using a specialized Scala-based Domain Specific
Language (Ocular/Joern dialect) to track data flow and identify security flaws.

\textit{Weaknesses and Limitations:} The primary weakness of Joern,
when considered for purely architectural dependency extraction, is
its sheer weight and complexity. Generating CFGs and PDGs requires a
deep understanding of the runtime execution semantics of the
language, which introduces massive computational overhead.
Furthermore, Joern does not use a single unified parser; it relies on
distinct, heavyweight frontends (e.g., Eclipse JDT for Java, Eclipse
CDT for C/C++) to generate the CPG for different languages. This
violates the goal of maintaining a lightweight, truly
language-agnostic extraction pipeline.

\section{Direct Precursors: Skullian and Entangle}

This thesis represents the culmination and third iteration of a
continuous research effort initiated by Francesco Refolli and later
advanced by Luca Teruzzi, building upon two previous prototypes
which laid the critical foundation for the current solution.

\textbf{Skullian} (Refolli, 2023) \cite{refolli_2023} was the initial
prototype of this research line, developed by Dr. Francesco Refolli
(who currently serves as co-advisor for this thesis). It attempted purely
language-agnostic dependency extraction by combining Tree-sitter for
parsing with Stack Graphs for name resolution.

\textit{How it works:} It parsed the source code with Tree-sitter and
used custom TSG rules to map the generated CST into a Stack Graph. It
then used the built-in Stack Graph path-finding algorithm to resolve
all references and build a dependency graph.

\textit{Weaknesses and Limitations:} The evaluation of Skullian
exposed severe practical and scalability limits of Stack Graphs for
architectural-level dependency extraction.
\begin{itemize}
  \item \textbf{Grammar Complexity}: Writing the language-specific
    TSG rules were highly verbose and complex, requiring thousands of
    lines of configuration for a single language like Java. This made
    the system difficult to maintain and extend.
  \item \textbf{Path Explosion and Cycle Detection Issues}: The Stack
    Graph size grew very big even for modest projects, causing
    reference resolution times to scale exponentially. To avoid
    infinite loops, the system relied on a cycle-detection heuristic.
    However, in real-world software, this heuristic frequently cut
    off valid resolution paths passing through the root node, leading
    to missed dependencies. Relaxing the path-similarity constraints
    (MAX\_SIMILAR\_PATH\_COUNT) inevitably resulted in system stalls or
    Out-Of-Memory (OOM) crashes, making large-scale projects unanalyzable.
  \item \textbf{Benchmark Results}:  comparison to Arcan, Skullian
    detected 22 out of 32 dependencies in the Pruijt et al.
    benchmark. On real-world software (JUnit4, JUnit5, and Antlr4),
    the resulting dependency graphs exhibited low structural
    similarity to Arcan's baseline, with Jaccard similarity indexes
    ranging between ~20\% and 48\% (specifically, 19.73\% for JUnit5 and
    25.42\% for JUnit4), making the approach inefficient and
    impractical for complex architectural analysis
\end{itemize}

\textbf{Entangle} (Teruzzi, 2024) \cite{teruzzi_2024} represents the
second iteration of this research line, developed by Luca Teruzzi
with the explicit goal of overcoming the critical performance and
memory scalability barriers encountered by Skullian.

\textit{How it works:} To achieve this, Entangle completely abandoned
Stack Graphs, introducing a custom, highly efficient two-phase name
resolution model.
\begin{itemize}
  \item In the first phase, it performs a static traversal of the
    source code's CST (generated via Tree-sitter) to reconstruct a
    hierarchical \textbf{Component Graph} representing packages,
    classes, and methods, while storing unresolved references as
    nested, unresolved dependency queries (modeled via recursive
    QueryData structures). A stack-based symbol table is maintained
    during this phase to map local variables to their respective types.
  \item \textbf{Phase Two: Scope Climbing Resolution}: In the second
    phase, the engine resolves these queries directly on the
    component graph using a custom \texttt{Lookup and Lookdown}
    algorithm. The lookup step ascends the component hierarchy to
    resolve the root of the query (handling imports as direct
    architectural bridges), while the lookdown step descends through
    the target component's children and inheritance chains to resolve
    nested members (e.g., method calls and field accesses).
\end{itemize}

To remain strictly language-agnostic, Entangle relies on a
configuration \textbf{dispatch table} for each supported language.
Users explicitly map abstract architectural concepts (e.g. Class,
Method) to the concrete CST node identifiers defined by that
language's Tree-sitter grammar (e.g., class\_declaration or
function\_item), isolating language-specific differences in a
declarative configuration layer

\textit{Weaknesses and Limitations:} despite drastically improving
performance, by consistently outperforming Arcan by completing
analyses in approximately half the execution time on average,
Entangle still exhibits structural and functional limitations:

\begin{itemize}
  \item \textbf{Manual Dispatch Table Configuration}: The dispatch
    table requires manual configuration for each language, demanding
    deep knowledge of the specific Tree-sitter grammar and its node
    identifiers.
  \item \textbf{Lack of Formalization}: Entangle's internal model
    lacks a rigorous mathematical formalization, which can lead to
    inconsistencies in name resolution and dependency extraction.
  \item \textbf{No Support for Generics}: The prototype completely
    lacks the ability to resolve and propagate parameterized types
    (Generics), resulting in a notable accuracy loss on
    template-heavy codebases.
  \item \textbf{Monolithic Module Resolution}: The framework is
    currently unable to partition distinct package modules (such as
    multiple Rust crates) within a single analyzed directory, causing
    them to collapse into a single dependency tree and compromising
    namespace isolation.
\end{itemize}

\section{Positioning of the Proposed Work}

The framework presented in this thesis directly addresses the
limitations of both the broad state-of-the-art technologies and the
specific precursor prototypes.

By relying strictly on Tree-sitter and its cursor navigation, we
maintain a lightweight, single-frontend extraction process, avoiding
the heavy, language-specific pipelines and CFG computational overhead
of tools like Joern. Instead of facing the severe scalability issues
of Stack Graphs (as evidenced by Skullian), or resorting to lossy AST
translations (as seen in Curtis and Schiewe), we have formalized a
robust Intermediate Representation ($\mathcal{D}_{IR}$) based on
algebraic data types.

Crucially, compared to Entangle, we have completely eliminated the
need for a manual "dispatch table". Instead, our framework implements
a \textbf{Heuristic-based extraction}. Because Tree-sitter parsers
use highly predictable naming conventions across different grammars
(e.g., node names consistently containing \texttt{struct},
\texttt{class}, \texttt{function}, or \texttt{module}), we extract
constructs by applying heuristics directly on the original CST.
Coupled with a feature-based configuration to manage broad behavioral
differences (like identifying if a language uses directory-based
packages or inline modules), this makes the extraction phase truly
and natively language-agnostic, eliminating the manual configuration
burden for each new language.

Furthermore, this formalization inherits the exceptional execution
speed achieved by Entangle while systematically solving its major
structural flaws. The formal IR natively models complex,
multi-layered module hierarchies (solving the namespace collision
issues). The resolution engine has been rigorously expanded to
support generic types, evaluated access paths, and accurate type
bounds. Ultimately, this thesis elevates the experimental concepts of
previous iterations into a mathematically sound, provably decoupled,
and highly performant language-agnostic architectural analyzer.
